Câncer colorretal é um tipo de tumor que afeta o intestino grosso e o reto, partes do corpo responsáveis por processar resíduos alimentares e absorver de água e de nutrientes. Após avaliação dos sintomas, o diagnóstico é confirmado por exames como pesquisa de sangue oculto nas fezes, colonoscopia e biópsia. Para sucesso do tratamento, é fundamental o diagnóstico precoce da doença. No entanto, a ausência de sintomas nas fases iniciais, associada à falta de acesso a exames de rastreamento, contribui para o diagnóstico tardio.
O objetivo desse trabalho é o uso de redes neurais profundas para a detecção precoce do câncer colorretal, utilizando imagens de colonoscopia. A metodologia proposta envolve a comparação de different arquiteturas, desde as mais clássicas como AlexNet, como também o uso de residuais como ResNet. Ademais, o trabalho inclui o uso de camadas de convoluções quânticas para a classificação de câncer colorretal.
Os resultados obtidos comprovam como o uso do aprendizado de máquina pode ser útil para a detecção precoce do câncer colorretal, contribuindo para o diagnóstico e tratamentos mais adequados, aumentando as chances de cura e melhorando a qualidade de vida dos pacientes.